% AUTO-GENERATED. Edit the Markdown sources or notes.config.json.
% Sources: 01-Mathematics/2026-Probability-Theory/2026-09-24-TA-Session-1.md

\section{第一次习题课}

\subsection{四类抽样计数}

从 \(n\) 类对象中选取 \(k\) 个对象时，应先判断是否考虑顺序、是否允许重复：

\begin{longtable}[]{@{}llll@{}}
\toprule\noalign{}
抽样方式 & 是否有序 & 是否放回 & 方案数 \\
\midrule\noalign{}
\endhead
\bottomrule\noalign{}
\endlastfoot
Permutation & 是 & 否 & \(P(n,k)=\dfrac{n!}{(n-k)!}\) \\
Combination & 否 & 否 & \(\binom nk=\dfrac{n!}{k!(n-k)!}\) \\
Ordered sampling with replacement & 是 & 是 & \(n^k\) \\
Unordered sampling with replacement & 否 & 是 & \(\binom{n+k-1}{k}\) \\
\end{longtable}

最后一种情形等价于非负整数解计数：

\[
x_1+x_2+\cdots+x_n=k,
\qquad x_i\ge 0,
\]

其解的数量为

\[
\binom{n+k-1}{n-1}=\binom{n+k-1}{k}.
\]

\subsection{\texorpdfstring{重复元素排列：\texttt{aeemrry}}{重复元素排列：aeemrry}}

单词 \texttt{aeemrry} 共含 7 个字母，其中 \texttt{e} 和 \texttt{r} 各重复两次。

\subsubsection{全部不同排列}

如果暂时把重复字母视为不同对象，会得到 \(7!\) 种排列；交换两个 \texttt{e} 或两个 \texttt{r} 不产生新排列，因此

\[
N=\frac{7!}{2!2!}=1260.
\]

\subsubsection{\texorpdfstring{包含连续子串 \texttt{eye}}{包含连续子串 eye}}

把 \texttt{eye} 看成一个整体，剩余对象为

\[
\{\text{eye},a,m,r,r\}.
\]

五个对象中两个 \texttt{r} 相同，所以

\[
N_{\text{eye}}=\frac{5!}{2!}=60.
\]

\subsubsection{\texorpdfstring{同时包含 \texttt{ram} 与 \texttt{eye}}{同时包含 ram 与 eye}}

把 \texttt{ram} 和 \texttt{eye} 分别看成一个整体，剩余一个 \texttt{r}。三个对象

\[
\{\text{ram},\text{eye},r\}
\]

互不相同，因此

\[
N_{\text{ram and eye}}=3!=6.
\]

\begin{quote}
\textbf{注：Block method}

要求某些元素连续出现时，可以先把它们压缩成一个 block，再对 blocks 与剩余元素排列。仍要检查不同 blocks 之间是否共享原始元素，以及剩余元素是否重复。
\end{quote}

\subsection{恰有一个空格的占位问题}

将 \(n\) 个有标号球独立、等概率地放入 \(n\) 个有标号格子。每个球有 \(n\) 种选择，因此总方案数为

\[
n^n.
\]

若恰有一个格子为空，则由于球数与格子数相同，必然出现以下占用结构：

\begin{itemize}
\tightlist
\item
  一个格子为空；
\item
  一个格子中有两个球；
\item
  其余 \(n-2\) 个格子各有一个球。
\end{itemize}

依次选择：

\begin{enumerate}
\def\labelenumi{\arabic{enumi}.}
\tightlist
\item
  空格子：\(n\) 种；
\item
  装两个球的格子：\(n-1\) 种；
\item
  同处一格的两个球：\(\binom n2\) 种；
\item
  将剩余 \(n-2\) 个球双射到剩余格子：\((n-2)!\) 种。
\end{enumerate}

故有利方案数为

\[
n(n-1)\binom n2(n-2)!
=\binom n2n!,
\]

所求概率为

\[
\boxed{
P=\frac{\binom n2n!}{n^n}
=\frac{(n-1)(n-1)!}{2n^{n-2}}
}.
\]

\begin{quote}
\textbf{警告：易错点}

只选择空格子还不够。必须同时确定哪个格子发生碰撞、哪两个球发生碰撞，以及其余球与格子的双射关系。
\end{quote}

\subsection{Stars and Bars：甜甜圈选择}

商店有 20 种甜甜圈，每种供应无限，要选择 12 个。令 \(x_i\) 表示第 \(i\) 种甜甜圈的数量，则

\[
x_1+x_2+\cdots+x_{20}=12,
\qquad x_i\ge 0.
\]

由 stars and bars，方案数为

\[
\boxed{\binom{20+12-1}{12}=\binom{31}{12}}.
\]

这里同一种甜甜圈不可区分，20 个种类可以区分，因此这是 unordered sampling with replacement。

\subsection{无标号分组}

\subsubsection{六人分成两个三人组}

先按顺序排列 6 人，再把相邻的每 3 人视为一组。组内顺序不重要，要除以 \((3!)^2\)；两个组没有标签，还要除以 \(2!\)：

\[
\frac{6!}{(3!)^2 2!}=10.
\]

\subsubsection{六人分成三个二人组}

同理，

\[
\frac{6!}{(2!)^3 3!}=15.
\]

\subsubsection{一般公式}

把 \(kn\) 个不同对象分成 \(n\) 个无标号组，每组 \(k\) 个，方案数为

\[
\boxed{
\frac{(kn)!}{(k!)^n n!}
}.
\]

分母中的 \((k!)^n\) 消除每组内部的排列，\(n!\) 消除组与组之间的排列。

\subsection{二项计数练习}

\subsubsection{十次抛硬币}

公平硬币独立抛掷 10 次，令 \(X\) 为正面次数，则

\[
X\sim\operatorname{Binomial}(10,1/2).
\]

恰有三次正面的概率为

\[
P(X=3)=\binom{10}{3}\left(\frac12\right)^{10}
=\frac{15}{128}
\approx0.1172.
\]

至多三次正面的概率为

\[
\begin{aligned}
P(X\le3)
&=\sum_{k=0}^{3}\binom{10}{k}\left(\frac12\right)^{10}\\
&=\frac{1+10+45+120}{1024}\\
&=\frac{11}{64}
=0.171875.
\end{aligned}
\]

\subsubsection{十人分成两支篮球队}

先选择一支 5 人队得到 \(\binom{10}{5}\)，但两队没有队名，交换两队不产生新划分，因此

\[
\frac{1}{2}\binom{10}{5}=126.
\]

\subsection{条件概率与 Bayes 公式}

\subsubsection{十年内从未中奖}

每周中奖概率为 \(\alpha=10^{-5}\)，各周相互独立。按每年 52 周计算，十年共有 520 次购票，因此

\[
P(\text{never win})=(1-10^{-5})^{520}\approx0.9948.
\]

单次中奖概率很小并不意味着长期中奖概率可以直接写成 \(520\alpha\)；后者只是 \(P(\text{at least one win})\) 的一阶近似。

\subsubsection{已知患流感，拥有狗的概率}

记 \(D\) 为拥有狗，\(F\) 为患流感。已知

\[
P(D)=0.4,
\qquad
P(F\mid D)=0.2,
\qquad
P(F)=0.15.
\]

由 Bayes 公式，

\[
P(D\mid F)
=\frac{P(F\mid D)P(D)}{P(F)}
=\frac{0.2\times0.4}{0.15}
=\frac{8}{15}
\approx0.5333.
\]

\subsubsection{两枚普通硬币与一枚双面正面硬币}

随机选择一枚硬币并抛掷。记 \(T\) 为选中双面正面硬币，\(H\) 为出现正面，则

\[
P(H)=P(H\mid T)P(T)+P(H\mid T^c)P(T^c)
=1\cdot\frac13+\frac12\cdot\frac23
=\frac23.
\]

观察到正面后，选中双面正面硬币的概率为

\[
P(T\mid H)
=\frac{P(H\mid T)P(T)}{P(H)}
=\frac{1/3}{2/3}
=\frac12.
\]

\subsection{两个证明型结论}

\subsubsection{在互斥分支上条件概率相同}

设 \(B_1,\ldots,B_k\) 两两互斥，\(B=\bigcup_{i=1}^kB_i\)，且

\[
P(A\mid B_i)=p,
\qquad P(B_i)>0.
\]

则

\[
\begin{aligned}
P(A\mid B)
&=\frac{P(A\cap B)}{P(B)}\\
&=\frac{\sum_{i=1}^kP(A\cap B_i)}{\sum_{i=1}^kP(B_i)}\\
&=\frac{\sum_{i=1}^kpP(B_i)}{\sum_{i=1}^kP(B_i)}\\
&=p.
\end{aligned}
\]

本质上，\(P(A\mid B)\) 是各个 \(P(A\mid B_i)\) 按 \(P(B_i\mid B)\) 加权得到的平均值；所有分支都等于 \(p\)，平均值仍为 \(p\)。

\subsubsection{条件独立推出无条件独立的一个充分条件}

设 \(C_1,\ldots,C_M\) 构成样本空间的 partition，并满足

\[
P(A\cap B\mid C_i)
=P(A\mid C_i)P(B\mid C_i),
\]

且 \(B\) 与每个 \(C_i\) 独立，即

\[
P(B\mid C_i)=P(B).
\]

使用全概率公式：

\[
\begin{aligned}
P(A\cap B)
&=\sum_{i=1}^{M}P(A\cap B\mid C_i)P(C_i)\\
&=\sum_{i=1}^{M}P(A\mid C_i)P(B\mid C_i)P(C_i)\\
&=P(B)\sum_{i=1}^{M}P(A\mid C_i)P(C_i)\\
&=P(A)P(B).
\end{aligned}
\]

因此 \(A\) 与 \(B\) 独立。

\begin{quote}
\textbf{警告：条件独立并不自动推出无条件独立}

上述证明依赖额外条件 \(B\perp C_i\)。如果不同分组同时影响 \(A\) 与 \(B\)，分组混合后可能产生相关性。
\end{quote}

\subsection{复习要点}

\begin{enumerate}
\def\labelenumi{\arabic{enumi}.}
\tightlist
\item
  计数题先判断对象是否可区分、结果是否有序、是否允许重复。
\item
  连续子串使用 block method；重复元素仍需除以内部不可区分排列数。
\item
  无标号分组要同时消除组内顺序和组间顺序。
\item
  占位问题要从 occupancy pattern 出发，列清楚空格、碰撞格、碰撞对象和剩余双射。
\item
  Bayes 公式的分母来自全概率公式；条件方向不能直接交换。
\item
  证明独立性最直接的方法是把目标化为 \(P(A\cap B)=P(A)P(B)\)。
\end{enumerate}

\subsection{术语表}

\begin{longtable}[]{@{}
  >{\raggedright\arraybackslash}p{(\linewidth - 4\tabcolsep) * \real{0.3333}}
  >{\raggedright\arraybackslash}p{(\linewidth - 4\tabcolsep) * \real{0.3333}}
  >{\raggedright\arraybackslash}p{(\linewidth - 4\tabcolsep) * \real{0.3333}}@{}}
\toprule\noalign{}
\begin{minipage}[b]{\linewidth}\raggedright
Term
\end{minipage} & \begin{minipage}[b]{\linewidth}\raggedright
中文
\end{minipage} & \begin{minipage}[b]{\linewidth}\raggedright
简要含义
\end{minipage} \\
\midrule\noalign{}
\endhead
\bottomrule\noalign{}
\endlastfoot
Block method & 捆绑法 & 把必须连续出现的对象视为一个整体进行排列 \\
Stars and bars & 隔板法 & 计算非负整数解或无序有放回抽样方案数 \\
Occupancy problem & 占位问题 & 将球分配到格子并研究各格占用数量 \\
Unlabeled groups & 无标号分组 & 组之间没有名称，交换整组不产生新方案 \\
Conditional independence & 条件独立 & 给定某一条件后，两个事件满足乘法分解 \\
Partition & 分割 & 两两互斥且并集为整个样本空间的一组事件 \\
\end{longtable}
